% !tp-example Basic usage
\documentclass{article}
\begin{document}
\wcTpOcamlTracer[src=/docs/components/ocaml-tracer/examples/recursion.ml]{}
\end{document}

% !tp-example Javascript API
\documentclass{article}
\begin{document}
\wcScript[type=module]|
import "/components/ocaml/ocaml.js";
import "/components/ocaml-tracer/ocaml-tracer.js";

const example = document.createDocumentFragment();

const element1 = document.createElement("tp-ocaml-tracer");
element1.setAttribute("src", "/docs/components/ocaml-tracer/examples/recursion.ml");
example.append(element1);

document.body.append(example);
|
\end{document}

% !tp-example Subexpressions on one line
\documentclass{article}
\begin{document}
\wcTpOcamlTracer[src=/docs/components/ocaml-tracer/examples/same-line.ml]{}
\end{document}

% !tp-example Shared arrays and references
\documentclass{article}
\begin{document}
\wcTpOcamlTracer[src=/docs/components/ocaml-tracer/examples/shared.ml]{}
\end{document}

% !tp-example Closures
\documentclass{article}
\begin{document}
\wcTpOcamlTracer[src=/docs/components/ocaml-tracer/examples/closures.ml]{}
\end{document}

% !tp-example Pattern matching
\documentclass{article}
\begin{document}
\wcTpOcamlTracer[src=/docs/components/ocaml-tracer/examples/patterns.ml]{}
\end{document}

% !tp-example Caught exception
\documentclass{article}
\begin{document}
\wcTpOcamlTracer[src=/docs/components/ocaml-tracer/examples/caught.ml]{}
\end{document}

% !tp-example Uncaught exception
\documentclass{article}
\begin{document}
\wcTpOcamlTracer[src=/docs/components/ocaml-tracer/examples/error.ml]{}
\end{document}

% !tp-example Replay a saved trace
\documentclass{article}
\begin{document}
\wcTpOcamlTracer[trace=/docs/components/ocaml-tracer/examples/recursion.trace.json]{}
\end{document}

% !tp-example Nested calls and a separate query
\documentclass{article}
\begin{document}
Record execution, then choose Start at query to skip the warmup calculation and reach fact(4), whose result is 24. Use Next to enter a nested call: Step over completes that call without visiting its body, while Step out returns from the active function to its caller. Use Previous to revisit the skipped steps; First also restores access to initialization. All steps remain available after saving and reloading the trace.


\wcTpOcamlTracer[src=/docs/components/ocaml-tracer/examples/nested-calls.ml,query=print_int (fact 4);;]{}
\end{document}

% !tp-example Constant expressions and intermediate values
\documentclass{article}
\begin{document}
Record execution and use Next to follow the highlighted operations and their Expression value until the result is 57. Previous revisits each calculation; Save trace and Load trace preserve the intermediate values without running the program again. Follow the order reported by the OCaml engine: it differs from the JavaScript and Scheme examples.


\wcTpOcamlTracer[src=/docs/components/ocaml-tracer/examples/constant-expressions.ml]{}
\end{document}
